% from ega0-1-fr.tex
\oldpage[0\textsubscript{I}]{38}
généralement des $\mathcal{C}$-Modules, si $\mathcal{C}$ est le centre de
$\mathcal{A}$). La fibre
$(\mathcal{F}\otimes_{\mathcal{A}}\mathcal{G})_x$ en tout point $x\in X$
s'identifie canoniquement à
$\mathcal{F}_x\otimes_{\mathcal{A}_x}\mathcal{G}_x$ et on définit un
homomorphisme canonique fonctoriel
\[
  (\mathcal{H}\!om_{\mathcal{A}}(\mathcal{F},\mathcal{G}))_x
  \longrightarrow
  \operatorname{Hom}_{\mathcal{A}_x}(\mathcal{F}_x,\mathcal{G}_x)
\]
qui, en général, n'est ni injectif ni surjectif. Les bifoncteurs considérés
ci-dessus sont additifs et, en particulier, commutent aux sommes directes
finies~; $\mathcal{F}\otimes_{\mathcal{A}}\mathcal{G}$ est exact à droite en
$\mathcal{F}$ et en $\mathcal{G}$, commute aux limites inductives, et
$\mathcal{A}\otimes_{\mathcal{A}}\mathcal{G}$
(resp. $\mathcal{F}\otimes_{\mathcal{A}}\mathcal{A}$) s'identifie
canoniquement à $\mathcal{G}$ (resp. $\mathcal{F}$). Les foncteurs
$\mathcal{H}\!om_{\mathcal{A}}(\mathcal{F},\mathcal{G})$ et
$\operatorname{Hom}_{\mathcal{A}}(\mathcal{F},\mathcal{G})$ sont
\emph{exacts à gauche} en $\mathcal{F}$ et $\mathcal{G}$~; de façon précise,
si on a une suite exacte de la forme
$0\to\mathcal{G}'\to\mathcal{G}\to\mathcal{G}''$, la suite
\[
  0\longrightarrow
  \mathcal{H}\!om_{\mathcal{A}}(\mathcal{F},\mathcal{G}')
  \longrightarrow
  \mathcal{H}\!om_{\mathcal{A}}(\mathcal{F},\mathcal{G})
  \longrightarrow
  \mathcal{H}\!om_{\mathcal{A}}(\mathcal{F},\mathcal{G}'')
\]
est exacte, et si on a une suite exacte de la forme
$\mathcal{F}'\to\mathcal{F}\to\mathcal{F}''\to0$, la suite
\[
  0\longrightarrow
  \mathcal{H}\!om_{\mathcal{A}}(\mathcal{F}'',\mathcal{G})
  \longrightarrow
  \mathcal{H}\!om_{\mathcal{A}}(\mathcal{F},\mathcal{G})
  \longrightarrow
  \mathcal{H}\!om_{\mathcal{A}}(\mathcal{F}',\mathcal{G})
\]
est exacte, avec des propriétés analogues pour le foncteur
$\operatorname{Hom}$. En outre,
$\mathcal{H}\!om_{\mathcal{A}}(\mathcal{A},\mathcal{G})$ s'identifie
canoniquement à $\mathcal{G}$~; enfin, pour tout ouvert $U\subset X$, on a
\[
  \Gamma\bigl(U,\mathcal{H}\!om_{\mathcal{A}}(\mathcal{F},\mathcal{G})\bigr)
  =
  \operatorname{Hom}_{\mathcal{A}|U}(\mathcal{F}|U,\mathcal{G}|U).
\]

Pour tout $\mathcal{A}$-Module à gauche $\mathcal{F}$ (resp. à droite), on
appelle \emph{dual de $\mathcal{F}$} et on note $\check{\mathcal{F}}$ le
$\mathcal{A}$-Module à droite (resp. à gauche)
$\mathcal{H}\!om_{\mathcal{A}}(\mathcal{F},\mathcal{A})$.

Enfin, si $\mathcal{A}$ est un faisceau d'anneaux commutatifs,
$\mathcal{F}$ un $\mathcal{A}$-Module,
$U\mapsto\bigwedge^p\Gamma(U,\mathcal{F})$ est un préfaisceau dont le
faisceau associé est un $\mathcal{A}$-Module qui se note
$\bigwedge^p\mathcal{F}$ et s'appelle la \emph{puissance extérieure $p$-ème
de $\mathcal{F}$}~; on vérifie aisément que l'application canonique du
préfaisceau $U\mapsto\bigwedge^p\Gamma(U,\mathcal{F})$ dans le faisceau
associé $\bigwedge^p\mathcal{F}$ est \emph{injective}, et que pour tout
$x\in X$, on a
$(\bigwedge^p\mathcal{F})_x=\bigwedge^p(\mathcal{F}_x)$. Il est clair que
$\bigwedge^p\mathcal{F}$ est un foncteur covariant en $\mathcal{F}$.
\end{env}

\begin{env}[4.1.6]
\label{0.4.1.6-fr}
Supposons que $\mathcal{A}$ soit un faisceau d'anneaux non nécessairement
commutatifs, $\mathcal{J}$ un faisceau d'idéaux à gauche de $\mathcal{A}$,
$\mathcal{F}$ un $\mathcal{A}$-Module à gauche~; on note alors
$\mathcal{J}\mathcal{F}$ le sous-$\mathcal{A}$-Module de $\mathcal{F}$,
image de $\mathcal{J}\otimes_{\mathbf Z}\mathcal{F}$ (où $\mathbf Z$ est le
faisceau associé au préfaisceau constant $U\to\mathbf Z$) par l'application
canonique $\mathcal{J}\otimes_{\mathbf Z}\mathcal{F}\to\mathcal{F}$~; il est
clair que pour tout $x\in X$, on a
$(\mathcal{J}\mathcal{F})_x=\mathcal{J}_x\mathcal{F}_x$. Lorsque
$\mathcal{A}$ est commutatif, $\mathcal{J}\mathcal{F}$ est aussi l'image
canonique de $\mathcal{J}\otimes_{\mathcal{A}}\mathcal{F}\to\mathcal{F}$.
Il est immédiat que $\mathcal{J}\mathcal{F}$ est aussi le
$\mathcal{A}$-Module associé au préfaisceau
$U\to\Gamma(U,\mathcal{J})\Gamma(U,\mathcal{F})$. Si
$\mathcal{J}_1$, $\mathcal{J}_2$ sont deux faisceaux d'idéaux à gauche de
$\mathcal{A}$, on a
$\mathcal{J}_1(\mathcal{J}_2\mathcal{F})=(\mathcal{J}_1\mathcal{J}_2)\mathcal{F}$.
\end{env}

\begin{env}[4.1.7]
\label{0.4.1.7-fr}
Soit $(X_\lambda,\mathcal{A}_\lambda)_{\lambda\in L}$ une famille d'espaces
annelés~; pour tout couple $(\lambda,\mu)$, supposons donnée une partie ouverte
$V_{\lambda\mu}$ de $X_\lambda$, et un isomorphisme d'espaces annelés
\[
  \varphi_{\lambda\mu}:
  (V_{\mu\lambda},\mathcal{A}_\mu|V_{\mu\lambda})
  \xrightarrow{\sim}
  (V_{\lambda\mu},\mathcal{A}_\lambda|V_{\lambda\mu}),
\]
avec $V_{\lambda\lambda}=X_\lambda$, $\varphi_{\lambda\lambda}$ étant
l'identité. Supposons de plus que, pour tout triplet $(\lambda,\mu,\nu)$, si
on désigne par $\varphi'_{\mu\lambda}$ la restriction de
$\varphi_{\mu\lambda}$ à $V_{\lambda\mu}\cap V_{\lambda\nu}$,
$\varphi'_{\mu\lambda}$ soit un isomorphisme de
\[
  (V_{\lambda\mu}\cap V_{\lambda\nu},
   \mathcal{A}_\lambda|(V_{\lambda\mu}\cap V_{\lambda\nu}))
\]
sur
\[
  (V_{\mu\nu}\cap V_{\mu\lambda},
   \mathcal{A}_\mu|(V_{\mu\nu}\cap V_{\mu\lambda}))
\]
et que l'on ait
$\varphi'_{\lambda\nu}=\varphi'_{\lambda\mu}\circ\varphi'_{\mu\nu}$
(\emph{condition de recollement} pour les $\varphi_{\lambda\mu}$). On peut
tout d'abord considérer alors l'espace topologique obtenu par recollement (au
moyen des $\varphi_{\lambda\mu}$) des $X_\lambda$

